\documentclass[11pt]{article}

\usepackage[margin=1in]{geometry}
\usepackage[T1]{fontenc}
\usepackage[utf8]{inputenc}
\usepackage{lmodern}
\usepackage{amsmath,amssymb}
\usepackage{enumitem}
\usepackage{hyperref}

\hypersetup{colorlinks=true,linkcolor=blue,urlcolor=blue}

\title{Experimental Design}
\author{}
\date{}

\begin{document}
\maketitle

\section{Branches}

The experiment currently uses four branches:
\begin{itemize}[leftmargin=2em]
    \item \textbf{Static artificial}
    \item \textbf{Static campus\_area\_1}
    \item \textbf{Dynamic port}
    \item \textbf{Dynamic campus\_area\_2}
\end{itemize}

\section{Assignment Model}

Branches normally use one-to-one MAPF assignments when goals are \texttt{dispersed} or \texttt{clustered}. The project also supports a target-only \texttt{single} goal distribution mode, where all agents share one literal target cell and disappear there after arrival. What changes across branches is therefore both the \emph{positioning} of starts and goals and, when \texttt{single} is selected for targets, the assignment cardinality.

\section{Current Spatial Patterns}

\begin{itemize}[leftmargin=2em]
    \item \texttt{static\_artificial}: dispersed starts, dispersed goals.
    \item \texttt{static\_campus\_area\_1}: dispersed starts in one campus zone, clustered goals in the other campus zone.
    \item \texttt{dynamic\_port}: clustered starts, dispersed goals.
    \item \texttt{dynamic\_campus\_area\_2}: clustered starts in one campus zone, clustered goals in one different campus zone.
\end{itemize}

\section{Campus Semantics}

For the campus branches, zone colors are traversable and spawnable, white walkways are traversable but non-spawnable, gray regions are non-traversable, and the dynamic campus branch generates dynamic obstacles only inside the zone colors.

\section{Placement Rule}

Within a sampled run configuration, \texttt{dispersed} sets keep 8-neighbor separation internally. Clustered sets are controlled by the global \texttt{compact\_clustering} switch: when it is \texttt{True}, clustered sets are directly adjacent 8-neighbor-connected groups; when it is \texttt{False}, clustered sets stay connected while keeping one empty cell of separation in all 8 directions. The \texttt{single} mode applies only to targets and samples one shared target cell. Start-goal adjacency is allowed, but a start and its assigned goal cannot be the same cell.

\end{document}
